\section{Introduction}

Large Language Models (LLMs) are rapidly evolving toward agentic intelligence. Recent advances, such as GPT-5.2~\citep{gpt52report}, Claude Opus 4.5~\citep{opus45report}, Gemini 3 Pro~\citep{gemini3pro}, and Kimi K2-Thinking~\citep{moonshotai2025kimi}, demonstrate substantial progress in agentic capabilities, particularly in tool calling and reasoning. These models increasingly exhibit the ability to decompose complex problems into multi-step plans and to execute long sequences of interleaved reasoning and actions.


In this report, we introduce the training methods and evaluation results of Kimi K2.5. Concretely, we improve the training of K2.5 over previous models in the following two key aspects.




\textbf{Joint Optimization of Text and Vision.} A key insight from the practice of K2.5 is that joint optimization of text and vision enhances both modalities and avoids the conflict. Specifically, we devise a set of techniques for this purpose. During pre-training, in contrast to conventional approaches that add visual tokens to a text backbone at a late stage~\citep{bai2025qwen3vltechnicalreport,guo2025seed15vltechnicalreport}, we find early vision fusion with lower ratios tends to yield better results given the fixed total vision-text tokens. Therefore, K2.5 mixes text and vision tokens with a constant ratio throughout the entire training process. 


Architecturally, Kimi K2.5 employs MoonViT-3D, a native-resolution vision encoder incorporating the NaViT packing strategy~\citep{dehghani2023patchnpacknavit}, enabling variable-resolution image inputs. For video understanding, we introduce a lightweight 3D ViT compression mechanism: consecutive frames are grouped in fours, processed through the shared MoonViT encoder, and temporally averaged at the patch level. This design allows Kimi K2.5 to process videos up to 4 $\times$ longer within the same context window while maintaining complete weight sharing between image and video encoders.


During post-training, we introduce zero-vision SFT---text-only SFT alone activates visual reasoning and tool use. 
We find that adding human-designed visual trajectories at this stage hurts generalization. 
In contrast, text-only SFT performs better—likely because joint pretraining already establishes strong vision-text alignment, enabling capabilities to generalize naturally across modalities.
We then apply joint RL on both text and vision tasks. 
Crucially, we find visual RL enhances textual performance rather than degrading it, with improvements on MMLU-Pro and GPQA-Diamond. This bidirectional enhancement—text bootstraps vision, vision refines text—represents superior cross-modal alignment in joint training.









\textbf{Agent Swarm: Parallel Agent Orchestration.}
Most existing agentic models rely on sequential execution of tool calls. Even systems capable of hundreds of reasoning steps, such as Kimi K2-Thinking~\citep{moonshotai2025kimi}, suffer from linear scaling of inference time, leading to unacceptable latency and limiting task complexity. As agentic workloads grow in scope and heterogeneity—e.g., building a complex project that involves massive-scale research, design, and development—the sequential paradigm becomes increasingly inefficient.

To overcome the latency and scalability limits of sequential agent execution, Kimi K2.5 introduces \textit{Agent Swarm}, a dynamic framework for parallel agent orchestration. We propose a Parallel-Agent Reinforcement Learning (PARL) paradigm that departs from traditional agentic RL~\citep{moonshotai2025kimiresearcher}. In addition to optimizing tool execution via verifiable rewards, the model is equipped with interfaces for sub-agent creation and task delegation.  During training, sub-agents are frozen and their execution trajectories are excluded from the optimization objective; only the orchestrator is updated via reinforcement learning. This decoupling circumvents two challenges of end-to-end co-optimization: credit assignment ambiguity and training instability. Agent Swarm enables complex tasks to be decomposed into heterogeneous sub-problems executed concurrently by domain-specialized agents, transforming task complexity from linear scaling to parallel processing. In wide-search scenarios, Agent Swarm reduces inference latency by up to 4.5$\times$ while improving item-level F1 from 72.8\% to 79.0\% compared to single-agent baselines.

Kimi K2.5 represents a unified architecture for general-purpose agentic intelligence, integrating vision and language, thinking and instant modes, chats and agents. It achieves strong performance across a broad range of agentic and frontier benchmarks, including state-of-the-art results in visual-to-code generation (image/video-to-code) and real-world software engineering in our internal evaluations, while scaling both the diversity of specialized agents and the degree of parallelism. To accelerate community progress toward General Agentic Intelligence, we open-source our post-trained checkpoints of Kimi K2.5, enabling researchers and developers to explore, refine, and deploy scalable agentic intelligence.



